\NativeTitle{41}{期中综合素养评价（一）}{}{90}{5}{一 & 二 & 三 & 四 & 五}{ &  &  &  & }
\NativeSlot{278}{36}{%
\ExamGroup{一、填空。{\normalfont\kaishu\fontsize{10.8bp}{14bp}\selectfont （每空1分，共25分）}}

% p041_s01

\ExamQuestion{1.}{5}%
要使\(9\times(\square\div4)=72\)，方框里应该填\mbox{（\hspace{1.8em}）}。\par
}
\NativeSlot{333}{36}{%
% p041_s02

\ExamQuestion{2.}{5}%
同一平面内，两条直线的位置关系有\mbox{（\hspace{1.8em}）}和\mbox{（\hspace{1.8em}）}。当两条直线相交成直角时，这两条直线互相\mbox{（\hspace{1.8em}）}。\par
}
\NativeSlot{385.0}{59}{%
% p041_s03

\ExamQuestion{3.}{5}%
根据运算律或性质填空：\par\noindent （1）\(357+168+43+32=(357+\underline{\qquad})+(168+\underline{\qquad})\)。\par

% p041_s04

\ExamQuestion{（2）}{5}%
\(125\times(80+8)=125\times\underline{\qquad}+125\times\underline{\qquad}\)。\par

% p041_s05

\ExamQuestion{（3）}{5}%
\(a\times99+a=a\times(\underline{\qquad}+\underline{\qquad})\)。\par
}
\NativeSlot{485}{36}{%
% p041_s06

\ExamQuestion{4（1）}{5}%
一个平行四边形的两条邻边分别长6厘米和4厘米，它的周长是\mbox{（\hspace{1.8em}）}厘米。\par
}
\NativeSlot{538.5}{36}{%
% p041_s07

\ExamQuestion{4（2）}{5}%
一个等边三角形的周长是27厘米，它的边长是\mbox{（\hspace{1.8em}）}厘米。\par
}
\NativeSlot{565}{89}{%
% p041_s08

\ExamQuestion{5.}{5}%
在圆圈里填上“＞”“＜”或“＝”。\(125+25\times4\ \bigcirc\ (125+25)\times4\)；\allowbreak{}\(36\times9+36\times2\ \bigcirc\ 36\times9+2\)；\allowbreak{}\(45\times(12+8)\ \bigcirc\ 45\times12+45\times8\)；\allowbreak{}\(600\div(2\times5)\ \bigcirc\ 600\div2\times5\)。\par
}
\NativeSlot{718}{45}{%
% p041_s09

\ExamQuestion{6.}{5}%
一个直角三角形中，有一个锐角是\(45^\circ\)，另一个锐角是\mbox{（\hspace{1.8em}）}，按边分它是\mbox{（\hspace{1.8em}）}三角形。把这样两个完全一样的三角形拼成一个大三角形，大三角形的内角和是\mbox{（\hspace{1.8em}）}；拼成一个正方形，正方形的内角和是\mbox{（\hspace{1.8em}）}。\par
}
